% Ad Atticum 14.14 (ad-atticum-14-14) --- English translation
% Translated by Claude (Anthropic) from the Latin (Perseus).
% Style guide: STYLE.md.

\ciceroLetterOpener{CICERO ATTICO S. D.}

\ciceroSection{1} Tell me that same thing again. Our Quintus crowned at the Parilia! At the Parilia! Was he the only one? Even though you add Lamia --- which does indeed surprise me --- I want to know who the others were; though I am quite sure there can have been none but scoundrels. You will therefore explain this more carefully. As for me, by chance, when I had sent off a letter to you on the sixth before the Kalends at considerable length, about three hours later I received yours, and one of much weight indeed. So there is no need to write back that your jokes, full of wit, about the Vestorian school and the \textit{†Pherionian fashion†} at Puteoli made me laugh enough. Let us look at the more political matters \textit{[Greek: politik\=otera]}.

\ciceroSection{2} You defend the Bruti and Cassius as though I were finding fault with them --- whom I cannot praise enough. It is the faults of the situation I have collected, not of the men. For with the tyrant removed I see the tyranny still standing. The things he would not have done are being done; as about Clodius, in whose case I have it for certain that he not only would not have done it, but would not even have allowed it. Next will come Rufio of the Vestorian set, then Victor (whose name was never on the rolls), then the rest --- who not? We obey the docket-papers of one whose own slave we could not bear to be. For at the Liberalia, who could fail to come to the Senate? Suppose somehow he could have stayed away: when we did come, could we then have spoken our minds freely? Was it not in every way our duty to defend ourselves against the veterans who were standing by under arms, while we had no garrison of our own? That I did not approve of the sitting on the Capitol, you can witness. What then? Is the blame the Bruti's? Not in the least theirs --- but of other ``brutes'' who fancy themselves cautious and wise; for whom it was enough to rejoice; for some of them, even to add their congratulations; for none, to see it through.

\ciceroSection{3} But let us drop the past. Let us guard those men with every care and protection, and, as you advise, let us be content with the Ides of March --- which gave to our friends, those godlike men, a way of approach to heaven, but did not give liberty to the Roman People. Recall your own words. Do you not remember that you used to cry out that everything would be lost if he were given a public funeral? That was wise judgement. And so you see what has flowed from it.

\ciceroSection{4} You write that on the Kalends of June Antony will bring before the Senate the matter of the provinces, so that he himself shall hold both Gauls and the terms of both governorships be extended. Will a free decision be possible? If it is, I shall rejoice that liberty has been recovered; if it is not, what has that change of master brought me, beyond the joy I had with my eyes from the tyrant's just destruction?

\ciceroSection{5} You write that depredations are going on at the Temple of Ops; we too were seeing them at the time. So we have neither been liberated by those excellent men nor are we free. So the credit is theirs, the blame ours. And you urge me to write history, to assemble such vast crimes of those by whom even now we are beleaguered! Shall I be able to refrain from praising the very men who employed you as a witness to their seals? Nor, by Hercules, is it the petty cash that moves me, but to attack with reproach men of goodwill, of whatever sort they are, is a heavy thing.

\ciceroSection{6} But as for all my plans, as you write, I think we can settle them more surely by the Kalends of June. I shall be there for them, and shall strain every effort and exertion --- helped, of course, by your authority and influence and the perfect justice of the case --- to see that the decree of the Senate concerning the people of Buthrotum comes out as you describe. As for what you bid me think over, I shall indeed think it over, though in my previous letter I had given you something to think over. You, meanwhile, as though the Republic were already recovered, are restoring to your neighbours the Massiliotes what is theirs. These things can perhaps be restored by arms (how firm a hold we have on arms I do not know); by mere authority they cannot.

\ciceroSection{7} The short letter that was written afterwards by you was very welcome to me, about Brutus's letter to Antony and the same man's letter to you. Things look as though they could be better than they have been so far. But for ourselves we must look ahead --- where we are to be, and where we are now to take ourselves.
